\begin{Example}
\label{ex:stokes-dirichlet-complementing-condition}
考虑半空间 $z>0$ 中的 Stokes 方程组 $\Delta\bm{u}-\nabla p=0$, $\operatorname{div}\bm{u}=0$, 其中 $\bm{u}\in\mathbb{R}^3$, $p\in\mathbb{R}$, 并在 $z=0$ 上施加 Dirichlet 边界条件 $\bm{u}=0$. 如 Chapter~\ref{chap:characteristics} 中所示, 该方程组椭圆, 且阶数为 $6$. 现在假设有解 $\bm{u}(x,y,z)=\exp(i\xi_1x+i\xi_2y)\bm{v}(z)$, $p(x,y,z)=\exp(i\xi_1x+i\xi_2y)q(z)$, 其中当 $z\to\infty$ 时, $\bm{v}$ 和 $q$ 指数趋于零. 令 $\Sigma=(0,L_1)\times(0,L_2)$, 其中若 $\xi_i\ne0$, 则 $L_i$ 为 $2\pi/\xi_i$ 的整数倍, 若 $\xi_i=0$, 则任意选取. 得到
\begin{equation}\label{eq:stokes-complementing-periodic-energy-identity}
0=\int_{\Sigma\times\mathbb{R}^+}(\Delta\bm{u}-\nabla p)\cdot\bm{u}\,\mathrm{d}x\,\mathrm{d}y\,\mathrm{d}z=-\int_{\Sigma\times\mathbb{R}^+}|\nabla\bm{u}|^2\,\mathrm{d}x\,\mathrm{d}y\,\mathrm{d}z.
\end{equation}
这表明 $\bm{u}$ 为常数, 而与假设相容的唯一情形是 $\bm{u}=0$. 容易得到 $p$ 也为零. 因而具有 Dirichlet 边界条件的 Stokes 方程组满足补条件.
\end{Example}

\begin{Example}
\label{ex:biharmonic-noncomplementing-boundary-condition}
考虑半平面 $y>0$ 中的双调和方程 $\Delta\Delta u=0$, 并在直线 $y=0$ 上施加边界条件 $\Delta u=\frac{\partial}{\partial y}\Delta u=0$. 对每个 $\xi\in\mathbb{R}$, 函数 $u(x,y)=\exp(i\xi x-|\xi|y)$ 都是解. 因而补条件不成立.
\end{Example}

补条件成立与否并非总像前面的例子那样容易验证. 然而, 总可以将它归结为纯代数问题. 将待定形式 \eqref{eq:complementing-decaying-tangential-mode} 代入 \eqref{eq:complementing-frozen-principal-system}, 得到关于 $\bm{v}(\eta)$ 的 ODE 方程组, 可照常用待定形式 $\bm{v}(\eta)=\exp(\lambda\eta)\bm{v}_0$ 求解. 椭圆性意味着没有纯虚根 $\lambda$; 若方程组的系数为实数, 则正实部根与负实部根的数目必相等. 设 $\lambda_l^+(\bm{x}_0,\bm{\xi})$ 表示具有正实部的根. 于是得到 \eqref{eq:complementing-frozen-principal-system} 的 $m$ 个线性无关解, 形式为 $\exp(i\bm{\xi}\cdot(\bm{x}-\bm{x}_0)+\lambda_l^+\eta)\bm{u}_l$ (如通常那样, 若有重根, 可能需要引入 $\eta$ 的幂加以修改). 剩下需要检验的是, 这些解的某个线性组合是否满足边界条件, 这是纯代数问题. 补条件的等价代数刻画参见文献 \cite{ADN2}.
